% Generated draft: requires mathematical and source review.
\documentclass[11pt]{article}
\usepackage[margin=25mm]{geometry}
\usepackage{fontspec}
\setmainfont{Latin Modern Roman}
\usepackage{amsmath,amssymb,mathtools,unicode-math}
\setmathfont{Latin Modern Math}
\usepackage{enumitem,needspace,xurl,hyperref}
\setlength{\parindent}{0pt}
\setlength{\parskip}{5pt}
\setlength{\emergencystretch}{3em}
\newcommand{\N}{\mathbb{N}}
\newcommand{\Z}{\mathbb{Z}}
\newcommand{\Q}{\mathbb{Q}}
\newcommand{\R}{\mathbb{R}}
\newcommand{\C}{\mathbb{C}}
\newcommand{\F}{\mathbb{F}}
\newcommand{\A}{\mathbb{A}}
\newcommand{\PP}{\mathbb{P}}
\newcommand{\E}{\mathbb{E}}
\newcommand{\eps}{\varepsilon}
\newcommand{\dd}{\,\mathrm{d}}
\begin{document}
\section*{UnsolvedMath ID: 700869}
\textbf{Source problem code:} catalogue entry
\subsection{Definitions and mathematical statement}
For \(n\geq 1\), let
\[
\operatorname{SL}_n(\Z)=\{A\in M_n(\Z):\det A=1\},
\]
with matrix multiplication as the group operation. For matrices \(A_1,\ldots,A_m\) in this group, define
\[
\langle A_1,\ldots,A_m\rangle
=\{A_{i_1}^{\varepsilon_1}\cdots A_{i_k}^{\varepsilon_k}: k\geq 0,\ 1\leq i_j\leq m,\ \varepsilon_j\in\{-1,1\}\}.
\]
The empty product is the identity matrix, and the subgroup generated by an empty family is the trivial subgroup.

The problem asks for a single algorithm which, for every integer \(m\geq 0\) and every input
\[
(B,A_1,\ldots,A_m)\in\operatorname{SL}_3(\Z)^{m+1},
\]
given by the integer entries of these matrices, terminates and answers yes if and only if
\[
B\in\langle A_1,\ldots,A_m\rangle.
\]
Here \(B\) is the source's first matrix and the \(A_i\) are its remaining matrices. This distinguishes a target matrix explicitly, since an unordered finite set has no first element.

For context, the supplied source states that the analogous decision problem for \(\operatorname{SL}_4(\Z)\) is insoluble, giving as its reason that \(F_2\times F_2\) embeds into \(\operatorname{SL}_4(\Z)\), where \(F_2\) denotes a free group of rank two. This contextual assertion is retained as a supplied source claim, without independent verification. The registry labels the dimension-three problem open and the source calls it a well-known problem; these are supplied classifications.
\subsection{Short English statement}
Find an algorithm that always decides whether a specified matrix in \(\operatorname{SL}_3(\Z)\) belongs to the subgroup generated by finitely many other given matrices. Products may use generators and their inverses. The supplied source describes the corresponding dimension-four problem as undecidable.
\subsection{Sources}
\begin{itemize}
\item Local UnsolvedMath registry, permanent ID \(700869\), exact source problem code \texttt{catalogue entry}, title \emph{Kourovka Notebook Problem 12.50}. Supplied upstream revision: \texttt{d371ee917ac9647fbb06cd6fc529deec6f3dda82}. The record was updated on 2026-10-16 and labels the problem open; this label is supplied registry data, not an independently verified literature assessment.
\item \emph{The Kourovka Notebook (Unsolved Problems in Group Theory)}, original Problem 12.50, 12th issue (1992), proposed by G. S. Makanin. The supplied background also identifies the 21st issue edition (2026), \texttt{arXiv:1401.0300v48}. The supplied collection description names E. I. Khukhro and V. D. Mazurov as editors of the Issues 1--21 collection. These bibliographic details and the dimension-four comparison are supplied attribution; the referenced edition has not been independently consulted.
\end{itemize}
\end{document}
